% GENERATED FILE. Edit canonical lessons, then run npm run generate:books.
% canonical-source-hash: fnv1a64:02ac12fbfb3f8aeb
% canonical-lessons: IT-C123-rosa, IT-C123-viola, IT-C123-marrone, IT-C123-arancione, IT-C123-carino

\chapter{Pink, Purple, Brown, Orange, Pretty}
\label{ch:it-pink-purple-brown-orange-pretty}
\hlchaptermodality{123}

\begin{chapteropening}
\textbf{By the end of this chapter:} \emph{I can say pink, purple, brown, orange and pretty.}

Complete the last lesson of chapter 123: I can say pink, purple, brown, orange and pretty.
\end{chapteropening}

\section[rosa]{rosa — pink}

\label{lesson:IT-C123-rosa}



\begin{quote}
Before the new one: say the Italian for a cheek, then the Italian for lips.
\end{quote}

\subsection*{You'll want to know: rosa}
\textbf{rosa} — \textquotedblleft{}pink\textquotedblright{}.

\begin{cousinweb}
\textbf{rosa} $\leftarrow$ Latin \emph{rosa}, \textquotedblleft{}a rose\textquotedblright{} — inherited. (EtymDB 2.1) The colour is named after the flower.
\end{cousinweb}

\begin{grammarlens}[title={What you've built}]
One more word for this chapter.
\end{grammarlens}

\subsection*{Your turn}
\begin{itemize}
\raggedright
  \item \emph{Say it:} \emph{rosa}
  \item \emph{Say it:} \emph{rosa}, once more
  \item \emph{Say it:} say \emph{rosa}
  \item \emph{Recall:} say the Italian for a cheek, then the Italian for lips, then say \emph{rosa} again
\end{itemize}

\begin{tcolorbox}[breakable,colback=teal!4,colframe=teal!35!black,title={Before you move on}]
What does rosa mean? (\textquotedblleft{}Pink\textquotedblright{}.) Say it once more.
\end{tcolorbox}
\section[viola]{viola — purple}

\label{lesson:IT-C123-viola}



\begin{quote}
Before the new one: say the Italian for lips, then the Italian for pink.
\end{quote}

\subsection*{You'll want to know: viola}
\textbf{viola} — \textquotedblleft{}purple\textquotedblright{}.

\begin{cousinweb}
\textbf{viola} $\leftarrow$ Latin \emph{viola}, \textquotedblleft{}a violet\textquotedblright{} — inherited. (EtymDB 2.1) The colour is named after the flower.
\end{cousinweb}

\begin{grammarlens}[title={What you've built}]
One more word for this chapter.
\end{grammarlens}

\subsection*{Your turn}
\begin{itemize}
\raggedright
  \item \emph{Say it:} \emph{viola}
  \item \emph{Say it:} \emph{viola}, once more
  \item \emph{Say it:} say \emph{viola}
  \item \emph{Recall:} say the Italian for lips, then the Italian for pink, then say \emph{viola} again
\end{itemize}

\begin{tcolorbox}[breakable,colback=teal!4,colframe=teal!35!black,title={Before you move on}]
What does viola mean? (\textquotedblleft{}Purple\textquotedblright{}.) Say it once more.
\end{tcolorbox}
\section[marrone]{marrone — brown}

\label{lesson:IT-C123-marrone}



\begin{quote}
Before the new one: say the Italian for pink, then the Italian for purple.
\end{quote}

\subsection*{You'll want to know: marrone}
\textbf{marrone} — \textquotedblleft{}brown\textquotedblright{}.

\begin{cousinweb}
\textbf{marrone}, the colour, is named after the \emph{marrone}, a large sweet chestnut. Where the chestnut word came from is uncertain: EtymDB 2.1 records no origin for it, only that French \emph{marron} was borrowed from it.
\end{cousinweb}

\begin{grammarlens}[title={What you've built}]
One more word for this chapter.
\end{grammarlens}

\subsection*{Your turn}
\begin{itemize}
\raggedright
  \item \emph{Say it:} \emph{marrone}
  \item \emph{Say it:} \emph{marrone}, once more
  \item \emph{Say it:} say \emph{marrone}
  \item \emph{Recall:} say the Italian for pink, then the Italian for purple, then say \emph{marrone} again
\end{itemize}

\begin{tcolorbox}[breakable,colback=teal!4,colframe=teal!35!black,title={Before you move on}]
What does marrone mean? (\textquotedblleft{}Brown\textquotedblright{}.) Say it once more.
\end{tcolorbox}
\section[arancione]{arancione — orange (the colour)}

\label{lesson:IT-C123-arancione}



\begin{quote}
Before the new one: say the Italian for purple, then the Italian for brown.
\end{quote}

\subsection*{You'll want to know: arancione}
\textbf{arancione} — \textquotedblleft{}orange (the colour)\textquotedblright{}.

\begin{cousinweb}
\textbf{arancione} — built from \emph{arancia}, \textquotedblleft{}an orange\textquotedblright{}, with the ending \emph{-one}: the colour of the fruit. EtymDB 2.1 has no row for it; the parts are plain to see.
\end{cousinweb}

\begin{grammarlens}[title={What you've built}]
One more word for this chapter.
\end{grammarlens}

\subsection*{Your turn}
\begin{itemize}
\raggedright
  \item \emph{Say it:} \emph{arancione}
  \item \emph{Say it:} \emph{arancione}, once more
  \item \emph{Say it:} say \emph{arancione}
  \item \emph{Recall:} say the Italian for purple, then the Italian for brown, then say \emph{arancione} again
\end{itemize}

\begin{tcolorbox}[breakable,colback=teal!4,colframe=teal!35!black,title={Before you move on}]
What does arancione mean? (\textquotedblleft{}Orange (the colour)\textquotedblright{}.) Say it once more.
\end{tcolorbox}
\section[carino]{carino — pretty, cute}

\label{lesson:IT-C123-carino}



\begin{quote}
Before the new one: say the Italian for brown, then the Italian for orange.
\end{quote}

\subsection*{You'll want to know: carino}
\textbf{carino} — \textquotedblleft{}pretty, cute\textquotedblright{}.

\begin{cousinweb}
\textbf{carino} — built from \emph{caro}, \textquotedblleft{}dear\textquotedblright{}, with \emph{-ino}, \textquotedblleft{}little\textquotedblright{}: a little dear. EtymDB 2.1 has no row for its origin; the parts are plain to see.
\end{cousinweb}

\begin{grammarlens}[title={What you've built}]
One more word for this chapter.
\end{grammarlens}

\subsection*{Your turn}
\begin{itemize}
\raggedright
  \item \emph{Say it:} \emph{carino}
  \item \emph{Say it:} \emph{carino}, once more
  \item \emph{Say it:} say \emph{carino}
  \item \emph{Recall:} say the Italian for brown, then the Italian for orange, then say \emph{carino} again
\end{itemize}

\begin{tcolorbox}[breakable,colback=teal!4,colframe=teal!35!black,title={Before you move on}]
What does carino mean? (\textquotedblleft{}Pretty, cute\textquotedblright{}.) Say it once more.
\end{tcolorbox}
